\documentclass[a4paper,11pt]{jsarticle}

\usepackage[dvipdfmx]{graphicx}
\usepackage[dvipdfmx]{hyperref}
\usepackage{pxjahyper}
\usepackage{amsmath, amssymb}
\usepackage{booktabs}
\usepackage{url}
\usepackage{geometry}
\geometry{margin=25mm}

\title{EEG Decoding Survey and Competition Strategy}
\author{Daisuke Sano}
\date{\today}

\begin{document}

\maketitle

\section{Objective}

This paper reviews the relevant research and outlines model design strategies for tasks involving the classification of image categories based on EEG signals recorded during the presentation of visual stimuli.
In particular, taking into account the spatial and temporal structures unique to EEG, we summarise the pre-processing techniques, model architectures and learning strategies that are effective in improving classification accuracy.

*** Translated with www.DeepL.com/Translator (free version) ***

\section{今回の課題設定}

本課題では、画像刺激を見た際に記録されたEEG信号から、画像カテゴリを5クラスに分類する。

\begin{itemize}
    \item 入力：$17$ channels $\times$ $100$ time points
    \item 出力：5クラス分類
    \item 訓練データ：118,800 trials
    \item 検証データ：59,400 trials
    \item テストデータ：59,400 trials
    \item 評価指標：Top-1 accuracy
\end{itemize}

各サンプルは、1回の画像提示に対するEEG応答を表すと考えられる。

\section{EEGデータの特徴}

EEGは、頭皮上電極から記録される多チャネル時系列信号である。
そのため、画像とは異なり、空間方向の局所性よりも、チャネル間の関係と時間方向の波形変化を適切に扱う必要がある。

本課題の入力は以下のテンソルとして表される。

\begin{equation}
    X \in \mathbb{R}^{C \times T},
\end{equation}

ここで、$C=17$ はチャネル数、$T=100$ は時系列長である。

\section{データ解析}

まず、訓練データの基本統計量を確認した。
全体の平均、標準偏差、最小値、最大値は以下の通りであった。

\begin{table}[h]
\centering
\caption{訓練EEGデータの基本統計量}
\begin{tabular}{cccc}
\toprule
mean & std & min & max \\
\midrule
-0.052 & 0.812 & -159.0 & 44.0 \\
\bottomrule
\end{tabular}
\end{table}

分位点を確認したところ、99.99パーセンタイルはおよそ5であり、大部分の値は$[-5,5]$の範囲内に収まっていた。
一方で、最小値や最大値には極端な外れ値が存在した。

\section{前処理方針}

外れ値の影響を抑えるため、まず以下のclippingを検討する。

\begin{equation}
    X_{\mathrm{clip}} = \mathrm{clip}(X, -5, 5).
\end{equation}

また、各trial・各channelごとに時間方向で標準化するtrial-wise z-scoreも検討する。

\begin{equation}
    X'_{c,t} = \frac{X_{c,t} - \mu_c}{\sigma_c + \epsilon},
\end{equation}

ここで、$\mu_c$および$\sigma_c$は各trial内のchannel $c$ における時間方向の平均および標準偏差である。

\section{関連研究}

\subsection{Gifford et al., 2022}

今回の課題は、Gifford et al. による視覚刺激提示時のEEGデータセットをもとに構成されている。
このデータセットは、画像刺激に対するEEG応答を用いて、視覚情報がEEG信号にどの程度反映されるかを検証するために用いられている。

\subsection{EEGNet}

EEGNetは、EEG分類で広く使われる軽量な畳み込みニューラルネットワークである。
時間方向の畳み込みとチャネル方向のdepthwise convolutionを組み合わせることで、EEG信号の時間特徴と空間特徴を効率的に抽出する。

\subsection{Conformer}

Conformerは、畳み込み層とTransformerを組み合わせた時系列モデルである。
畳み込みにより局所的な時間特徴を抽出し、self-attentionにより長距離依存関係を捉えることができる。

EEGは時系列信号であるため、ConformerはEEG decodingにも有効である可能性がある。

\subsection{ARAYAのSpeech Decoding研究}

ARAYAによる非侵襲的speech decoding研究では、大規模EEGデータを用いて、EEG表現と音声表現を対応づけるモデルが提案されている。
この研究では、EEG信号に対してz-score、clipping、temporal convolution、Conformerなどが用いられており、本課題の前処理およびモデル設計にも参考になる。

\subsection{Algonauts Project 2023}

Algonauts Project 2023では、自然画像を見た際のfMRI応答を予測する課題が扱われた。
入力モダリティは異なるものの、視覚刺激と脳活動を対応づけるという点で、本課題と共通する。

\section{モデル設計方針}

本課題では、以下のモデルを段階的に検討する。

\begin{enumerate}
    \item Baseline 1D CNN
    \item Baseline + clipping
    \item Baseline + trial-wise z-score + clipping
    \item EEGNet
    \item TCN
    \item EEGNet + Conformer
    \item 複数seedおよび複数モデルのアンサンブル
\end{enumerate}

\section{実験結果}

実験結果は以下の表にまとめる。

\begin{table}[h]
\centering
\caption{実験結果}
\begin{tabular}{lccc}
\toprule
Model & Preprocess & Validation Accuracy & Test Accuracy \\
\midrule
Baseline & None & -- & -- \\
Baseline & Clip $[-5,5]$ & -- & -- \\
Baseline & Z-score + Clip & -- & -- \\
EEGNet & Clip $[-5,5]$ & -- & -- \\
EEGNet + Conformer & Clip $[-5,5]$ & -- & -- \\
\bottomrule
\end{tabular}
\end{table}

\section{考察}

まず、外れ値解析により、訓練データにはごく少数の極端な外れ値が存在することが分かった。
そのため、clippingはモデルの学習安定化に寄与する可能性がある。

また、EEGはチャネル方向と時間方向の構造を持つため、単純な1D CNNよりも、EEGNetやConformerのようにEEG信号の構造を明示的に利用するモデルが有効であると考えられる。

\section{今後の方針}

今後は、まずbaselineに対する前処理の効果を検証し、その後EEGNet、TCN、Conformerへ拡張する。
最終的には、異なるモデル構造および異なるseedで学習したモデルのsoftmax確率平均により、汎化性能の向上を狙う。

\begin{thebibliography}{9}

\bibitem{gifford2022}
Gifford, A. T., et al.
Human EEG recordings for object recognition.
\textit{NeuroImage}, 2022.

\bibitem{eegnet}
Lawhern, V. J., et al.
EEGNet: a compact convolutional neural network for EEG-based brain-computer interfaces.
\textit{Journal of Neural Engineering}, 2018.

\bibitem{conformer}
Gulati, A., et al.
Conformer: Convolution-augmented Transformer for Speech Recognition.
\textit{Interspeech}, 2020.

\bibitem{araya2024}
ARAYA Inc.
Scaling Law in Neural Data: Non-Invasive Speech Decoding with 175 Hours of EEG Data.
\textit{arXiv}, 2024.

\bibitem{algonauts2023}
The Algonauts Project 2023 Challenge.
\url{https://algonautsproject.com/2023/}

\end{thebibliography}

\end{document}